\documentclass[10pt,a4paper]{amsart}
\usepackage[margin=19mm]{geometry}
\usepackage{amsmath,amssymb,mathtools}
\usepackage[scr=rsfs,cal=euler]{mathalfa}
\usepackage{xfrac}
\usepackage[unicode,hidelinks,bookmarks=false]{hyperref}
\newcommand{\ie}{i.e.\ }
\newcommand{\cf}{cf.\ }
\newcommand{\resp}{respectively\ }
\newcommand{\Subsection}{\subsection}
\providecommand{\nobreakdash}{\nobreak\hbox{-}}
\setlength{\emergencystretch}{2em}
\multlinegap0pt
\usepackage{bm}
\usepackage{bbm}
\usepackage{dsfont}
\usepackage{xspace}
\usepackage{cancel}
\usepackage{mathtools}
\usepackage{amsthm}
\makeatletter
\@ifundefined{theorem}{\newtheorem{theorem}{Theorem}}{}
\@ifundefined{lemma}{\newtheorem{lemma}[theorem]{Lemma}}{}
\makeatother
\title[On the equivalence of Sobolev norms in infinite dimensions]{On the equivalence of Sobolev norms in infinite dimensions}
\author{Zhouzhe Wang}
\date{Source: arXiv:2608.24702v1 (CC BY 4.0)}
\begin{document}
\maketitle

\noindent\emph{Source and licence.} On the equivalence of Sobolev norms in infinite dimensions. Version arXiv:2608.24702v1 (CC BY 4.0), distributed under the Creative Commons Attribution 4.0 International licence, as recorded in the arXiv licence metadata for this exact version and shown on the abstract page https://arxiv.org/abs/2608.24702v1. Full rights notice: licenses/arxiv-2608.24702-CC-BY-4.0.txt.\par\smallskip

\noindent\textit{Editorial scope.} Counted unit: the Theorem whose source label is \texttt{thm:vector-malliavin-equivalence-all-p-k} in the article (source0.tex lines 1185--1202, source label \texttt{thm:vector-malliavin-equivalence-all-p-k}), delivered together with the article's own complete original proof of it (source0.tex lines 1204--1310, one \texttt{proof} environment of 107 lines). No mathematics is written, repaired or re-derived: statement and proof are the article's own text line for line. Required context retained uncounted: thm:scalar-malliavin-equivalence-all-p-k (lines 997--1008); lem:gaussian-randomization-hilbert (lines 1107--1127). Every definition, display and label of those lines is retained unchanged. Omitted from this package and not redistributed: 1--996, 1009--1106, 1128--1184, 1203--1203, 1311--1641. Nothing inside a retained excerpt is omitted or altered: every retained body is delivered exactly as the article's lines are written, so no omission receipt of any kind accompanies this package. Citations: the retained excerpts carry no citation command at all, so no bibliography entry of the article is transcribed and no reference is redistributed. Reference adaptation: none was needed and none was made; every reference inside the delivered unit resolves inside it, so no unresolved reference or citation is produced. Printed slips kept literal: no printed word, symbol, hypothesis or number of the counted statement or of its proof was altered. Layout and class: the article's own class is \texttt{article} with packages including geometry, amsmath, amssymb, amsfonts, amsthm, mathtools, mathrsfs, xcolor, microtype, hyperref; the wrapper is the lane's pinned \texttt{amsart} editorial wrapper at 10pt on A4, which restates the theorem-like environments the retained text needs and adds the article's own macro definitions that the retained text needs (none), with extra wrapper packages (bm, bbm, dsfont, xspace, cancel, mathtools). The delivered numbering is the unit's own. Assets: no retained span contains a figure environment, a graphics-inclusion command, a PDF-page inclusion, a table or a TikZ picture; no image file is copied into this package, the package declares no asset dependency and no image file is redistributed with it. Licence: version arXiv:2608.24702v1 of this article is distributed under the Creative Commons Attribution 4.0 International licence, as recorded in the arXiv licence metadata for this exact version and shown on the abstract page https://arxiv.org/abs/2608.24702v1. A full rights notice is delivered at licenses/arxiv-2608.24702-CC-BY-4.0.txt. Characters: every retained excerpt is pure ASCII, so the delivered engine, XeLaTeX with its default Latin Modern text font, reports no missing character.
\begin{theorem}[Scalar Malliavin--Sobolev norm equivalence]
	\label{thm:scalar-malliavin-equivalence-all-p-k}
	Let $p\in[1,+\infty)$ and $k\geqslant2$.  Then, for every real separable Hilbert space $\mathfrak{H}$, every isonormal Gaussian process $W$ over $\mathfrak{H}$ and every $F\in\mathcal{S}$,
	$$
	\|F\|_{\mathcal{G}(k,p)}
	\leqslant
	\|F\|_{\mathcal{D}(k,p)}
	\leqslant
	(1+L_{p,k})\|F\|_{\mathcal{G}(k,p)}.
	$$
	In particular, the equivalence constant depends only on $p$ and $k$ and is independent of $\dim\mathfrak{H}$.
\end{theorem}

\begin{lemma}
	\label{lem:gaussian-randomization-hilbert}
There exist constants
	$0<\alpha_p\leqslant\beta_p<+\infty$, depending only on $p$, such that
	for every real Hilbert space $U$, every $M\in\mathbb{N}$, and every
	$u_1,\ldots,u_M\in U$,
	\[
	\alpha_p
	\left(\sum_{\ell=1}^{M}\|u_\ell\|_U^2\right)^{\frac12}
	\leqslant
	\left(
	\int_{\mathbb{R}^M}
	\left\|\sum_{\ell=1}^{M}y_\ell u_\ell\right\|_U^p
	\,\mathrm{d}\gamma_M(\mathbf{y})
	\right)^{\frac1p}
	\leqslant
	\beta_p
	\left(\sum_{\ell=1}^{M}\|u_\ell\|_U^2\right)^{\frac12}.
	\]
	For $p=2$, one may take $\alpha_2=\beta_2=1$.
\end{lemma}

\begin{theorem}[Hilbert-valued Malliavin--Sobolev norm equivalence]
	\label{thm:vector-malliavin-equivalence-all-p-k}
	Let $p\in[1,+\infty)$ and $k\geqslant2$.  Let $\mathfrak{H}$ and $V$ be arbitrary real separable Hilbert spaces and let $W$ be an isonormal Gaussian process over $\mathfrak{H}$.  Then, for every $F\in\mathcal{S}_{V}$,
	$$
	\|F\|_{\mathcal{G}(k,p)(V)}
	\leqslant
	\|F\|_{\mathcal{D}(k,p)(V)}
	\leqslant
	K_{p,k}\|F\|_{\mathcal{G}(k,p)(V)},
	$$
	where
	$$
	K_{p,k}
	\triangleq
	1+\frac{\beta_{p}}{\alpha_{p}}L_{p,k}.
	$$
	In particular, $K_{p,k}$ depends only on $p$ and $k$ and is independent of $\dim\mathfrak{H}$ and $\dim V$.
\end{theorem}

\begin{proof}
	The first inequality is immediate.  We prove the reverse one.
	
	Let $F\in\mathcal{S}_{V}$.  Since $F$ has finite-dimensional range, after replacing the original spanning family in $V$ by an orthonormal basis of its span, we may write
	$$
	F=\sum_{\ell=1}^{M}F_{\ell}v_{\ell},
	$$
	where $v_{1},\ldots,v_{M}$ are orthonormal in $V$ and $F_{1},\ldots,F_{M}\in\mathcal{S}$.
	
	For $\textbf{y}=(y_{1},\ldots,y_{M})\in\mathbb{R}^{M}$ define the scalar-valued smooth random variable
	$$
	F_{\textbf{y}}
	\triangleq
	\sum_{\ell=1}^{M}y_{\ell}F_{\ell}.
	$$
	For every $j\geqslant1$,
	$$
	D^{j}F_{\textbf{y}}
	=
	\sum_{\ell=1}^{M}y_{\ell}D^{j}F_{\ell}.
	$$
	Fix $j\in\{1,\ldots,k-1\}$.  By the scalar estimate in the proof of Theorem~\ref{thm:scalar-malliavin-equivalence-all-p-k}, for every $\textbf{y}\in\mathbb{R}^{M}$,
	$$
	\|D^{j}F_{\textbf{y}}\|_{L^{p}(\Xi;\mathfrak{H}^{\otimes j})}
	\leqslant
	L_{p,k,j}
	\left[
	\|F_{\textbf{y}}\|_{L^{p}(\Xi)}
	+
	\|D^{k}F_{\textbf{y}}\|_{L^{p}(\Xi;\mathfrak{H}^{\otimes k})}
	\right].
	$$
	Taking the $L^{p}(\mathbb{R}^{M},\gamma_{M})$ norm in the parameter $\textbf{y}$ and using Minkowski's inequality, we get
	$$
	\begin{aligned}
		&\left(
		\int_{\mathbb{R}^{M}}
		\|D^{j}F_{\textbf{y}}\|_{L^{p}(\Xi;\mathfrak{H}^{\otimes j})}^{p}
		\,\mathrm{d}\gamma_{M}(\textbf{y})
		\right)^{\frac1p}
		\\
		&\quad\leqslant
		L_{p,k,j}
		\left(
		\int_{\mathbb{R}^{M}}
		\|F_{\textbf{y}}\|_{L^{p}(\Xi)}^{p}
		\,\mathrm{d}\gamma_{M}(\textbf{y})
		\right)^{\frac1p}
		\\
		&\qquad\quad+
		L_{p,k,j}
		\left(
		\int_{\mathbb{R}^{M}}
		\|D^{k}F_{\textbf{y}}\|_{L^{p}(\Xi;\mathfrak{H}^{\otimes k})}^{p}
		\,\mathrm{d}\gamma_{M}(\textbf{y})
		\right)^{\frac1p}.
	\end{aligned}
	$$
	By Fubini's theorem and Lemma~\ref{lem:gaussian-randomization-hilbert}, applied pointwise in $\omega\in\Xi$, the left-hand side is bounded from below by
	$$
	\alpha_{p}
	\|D^{j}F\|_{L^{p}(\Xi;\mathfrak{H}^{\otimes j}\otimes V)}.
	$$
	Indeed, since $v_{1},\ldots,v_{M}$ are orthonormal,
	$$
	\sum_{\ell=1}^{M}
	\|D^{j}F_{\ell}(\omega)\|_{\mathfrak{H}^{\otimes j}}^{2}
	=
	\|D^{j}F(\omega)\|_{\mathfrak{H}^{\otimes j}\otimes V}^{2}.
	$$
	Similarly, the two terms on the right-hand side are bounded from above by
	$$
	\beta_{p}\|F\|_{L^{p}(\Xi;V)}
	$$
	and
	$$
	\beta_{p}
	\|D^{k}F\|_{L^{p}(\Xi;\mathfrak{H}^{\otimes k}\otimes V)},
	$$
	respectively.  Hence
	$$
	\|D^{j}F\|_{L^{p}(\Xi;\mathfrak{H}^{\otimes j}\otimes V)}
	\leqslant
	\frac{\beta_{p}}{\alpha_{p}}L_{p,k,j}
	\left[
	\|F\|_{L^{p}(\Xi;V)}
	+
	\|D^{k}F\|_{L^{p}(\Xi;\mathfrak{H}^{\otimes k}\otimes V)}
	\right].
	$$
	Summing this estimate over $j=1,\ldots,k-1$ gives
	$$
	\sum_{j=1}^{k-1}
	\|D^{j}F\|_{L^{p}(\Xi;\mathfrak{H}^{\otimes j}\otimes V)}
	\leqslant
	\frac{\beta_{p}}{\alpha_{p}}L_{p,k}
	\|F\|_{\mathcal{G}(k,p)(V)}.
	$$
	Adding the zeroth- and $k$-th-order terms yields
	$$
	\|F\|_{\mathcal{D}(k,p)(V)}
	\leqslant
	\left(1+\frac{\beta_{p}}{\alpha_{p}}L_{p,k}\right)
	\|F\|_{\mathcal{G}(k,p)(V)},
	$$
	which proves the theorem.
\end{proof}

\end{document}
